\documentclass[a4paper,12pt]{article}
\usepackage{docsTemplate}
\usepackage{longtable}

\setTitle{Verbale Interno 27/03/2026}
\setAuthors{Nenad Radulovic}
\setVerificators{Bianca Zaghetto}
\setApprovation{Filippo Compagno}
\setVersion{1.0}
\setType{Interno}
\setPlace{Discord}
\setTime{14.00}
\setDestination{Team Three Technologies \\ & Prof. Tullio Vardanega \\ & Prof. Riccardo Cardin}
\setAbsents{Nessuno}

\begin{document}
    \include{memoTitlepage}

    \newpage
    \tableofcontents
    \newpage

    \section{Ordine del giorno}
        \begin{itemize}
            \item Richiesta riunione con il prof. Cardin
            \item Miglioramenti alla documentazione 
            \item Cadenza degli incontri 
            \item Progettazione del renderer Angular
            \item Inizio stesura Specifica Tecnica
        \end{itemize}

    \section{Svolgimento}
        \begin{longtable}{P{5cm}|P{10cm}} 
            \textbf{Richiesta riunione con il prof. Cardin} & Si è deciso di richiedere al professor Cardin un 
            colloquio per chiarire dei dubbi che sono sorti durante la prima fase di progettazione.
            È stato incentivato l'uso del canale Telegram apposito per segnalare ulteriori dubbi qualora
            sorgessero nel periodo di attesa al colloquio.\\\\

            \textbf{Miglioramenti alla documentazione} & Basandosi sulla correzione del RTB si è deciso di apportare una serie di modifiche ai 
            documenti e al metodo di lavoro:
            \begin{itemize}
                \item Aggiungere un riferimento preciso ai verbali esterni e interni nei punti in cui non viene 
                    indicato con precisione a quale dei verbali ci si riferisce.
                \item Rimuovere il registro delle modifiche per i prossimi verbali
                \item Cambio della logica di esposizione della vista web del repo documentale, da concentrarsi
                    solo sulla baseline più recente di progetto.
                \item Aggiunta di un versionamento più rigido dei template dei verbali e delle presentazioni tramite
                    l'aggiunta degli stessi alla repository Github dedicata alla documentazione.
            \end{itemize}\\\\

            \textbf{Cadenza degli incontri} & Accogliendo i suggerimenti del professor Vardanega, durante la presentazione
            del RTB, sono stati decisi i seguenti provvedimenti:
            \begin{itemize}
                \item Fissare settimanalmente una riunione interna con tutto il gruppo per stabilire gli obbiettivi per 
                    il periodo a venire e per discutere le questioni che non si è riusciti a risolvere durante il periodo
                    trascorso. L'incontro è stato fissato per ogni venerdì alle 14:00.
                \item Il gruppo ha deciso di ritrovarsi quotidianamente, due volte al giorno per lavorare in sincrono sulle
                    attività, in modo che i membri possano supportasi e imparare tra di loro. Per venire incontro alle 
                    esigenze dei membri si è deciso di programmare due incotri, il primo alle 09:00 e il secondo alle 14:00, in modo da permettere a ogni membro di partecipare ad almeno una "sessione di lavoro" al giorno. 
                    I membri del gruppo si impegnano a partecipare alle "sessioni di lavoro" per almeno tre ore, in modo da riuscire ad ottenere dei progressi tangibili per ogni sessione.
            \end{itemize}\\\\

            \textbf{Progettazione del renderer Angular} & Alcuni membri del gruppo hanno preso in carico il compito di iniziare la progettazione
                        del renderer Angular. Si è resa nota la necessità di approfondire
                        al meglio le soluzioni offerte dal framework per implementare lo use case di visualizzazione dell'anteprima 
                        di documenti e allegati.\\\\
            \textbf{Inizio stesura Specifica Tecnica} & Si è ritenuto necessario iniziare con la stesura del documento
            delle Specifica Tecnica, in quanto i membri del gruppo si sono suddivisi in due sottogruppi per lavorare
            in parallelo ad aspetti diversi e si vuole avere fin dall'inizio un tracciamento delle scelte tecniche fatte, in modo da permettere a qualsiasi componente del gruppo di avere una chiara visione anche degli aspetti del progetto a cui non si è dedicato personalmente.\\\\
        \end{longtable}


    \section{Decisioni}
        \rowcolors{2}{lightgray}{white}
        \begin{tabularx}{\textwidth}{|l|X|l|}
            \hline
            \textbf{Identificativo} & \textbf{Descrizione}\\
            \hline
            VI\_20260327.d1 & Richiedere un colloquio con il professor Cardin \\
            \hline
            VI\_20260327.d2 & Rimuovere dai verbali il registro delle modifiche  \\
            \hline
            VI\_20260327.d3 & Riunione interna a cadenza settimanale il venerdì alle 14:00\\
            \hline
            VI\_20260327.d4 & Sessione di lavoro giornaliera alle 09:00 e alle 14:00 \\
            \hline
            VI\_20260327.d5 & Inizio progettazione renderer Angular\\
            \hline
            VI\_20260327.d6 & Inizio stesura Specifica Tecnica\\
            \hline
        \end{tabularx}

    \section{Attività}
        \rowcolors{2}{lightgray}{white}
        \begin{tabularx}{\textwidth}{|>{\raggedright\arraybackslash}p{4cm}|X|l|}
            \hline
            \textbf{Incaricato} & \textbf{Task Todo} & \textbf{Identificativo} \\
            \hline
            \textbf{Tutti} & Correggere i riferimenti ai verbali esterni/interni & VI\_20260327.a1 \\
            \hline
            \textbf{Filippo Compagno} & Cambio logica esposizione del repo documentale & VI\_20260327.a2 \\
            \hline
            \textbf{Francesco Balestro, Sara Gioia Fichera, Andrea Masiero, Mattia Olivia Medin} & Progettazione renderer Angular& VI\_20260327.a3 \\
            \hline
            \textbf{Filippo Compagno, Nenad Radulovic} & Inizio stesura Specifica Tecnica & VI\_20260327.a4 \\
            \hline
        \end{tabularx}
\end{document}
